%% ====================================================================
\section{Error Analysis}
\label{sec:errors}
%% ====================================================================

\subsection{Truncation Error (Finite $N$)}
\label{sec:trunc}

The dominant systematic error arises from truncating the Dirichlet series at $N$ terms.

\paragraph{Without Euler--Maclaurin correction.}
\begin{equation}\label{eq:trunc-no-em}
  |\eta_N(s) - \eta(s)| \le C(s)\,(N{+}1)^{-\beta},
\end{equation}
where $C(s) = O(1)$ for bounded $|s|$. For $\beta = 0.5$ and $N = 10^4$: truncation
error $\sim 10^{-2}$.

\paragraph{With Euler--Maclaurin correction ($p$ terms).}
\begin{equation}\label{eq:trunc-em}
  |\eta_N^{\mathrm{EM}}(s) - \eta(s)|
  \le C_p\,|s|^{2p-1}\,N^{-(\beta + 2p - 1)},
\end{equation}
where the factor $|s|^{2p-1}$ accounts for the rising factorial in the correction terms
(see Section~\ref{sec:dirichlet}). For $\beta = 0.5$, $p = 5$, $N = 10^4$, and
$t = 100$: truncation error $\sim 10^{9} \times 10^{-38} = 10^{-29}$. See the precision
table in Section~\ref{sec:dirichlet} (Table~\ref{tab:em-precision}) for additional
parameter combinations.

This error is purely systematic and can be reduced by increasing $N$ (at the cost of
more computation) or by choosing $p$ optimally for the given $(N, t)$ pair.


\subsection{Floating-Point Precision}
\label{sec:trig}

IEEE~754 double-precision arithmetic provides approximately $15.9$ significant decimal
digits ($53$-bit mantissa). The key sources of floating-point error are:

\paragraph{Summation error.}
Accumulating $N$ terms in floating point introduces rounding error of order
$N\,\eps_{\mathrm{mach}} \sim N \times 10^{-16}$. For $N = 10^4$: rounding error
$\sim 10^{-12}$. For $N = 10^6$: rounding error $\sim 10^{-10}$. This can be mitigated
by Kahan compensated summation~\cite{kahan1965}, which reduces the error to
$O(\eps_{\mathrm{mach}})$ independent of~$N$.

\paragraph{Trigonometric evaluation and the FP64 precision wall.}
$\cos(t\ln n)$ and $\sin(t\ln n)$ are evaluated by the standard library, with error
$O(\eps_{\mathrm{mach}})$ per evaluation. For large arguments ($t\ln n \gg 1$), the
argument reduction step introduces additional error proportional to
$t\ln(n)\,\eps_{\mathrm{mach}}$. This constitutes a fundamental precision wall for
double-precision computation:
\begin{center}
\begin{tabular}{@{}rll@{}}
\toprule
$t$ & Argument $t\ln n$ (at $n = 10^4$) & Effective digits lost \\
\midrule
$10^3$ & $\sim 10^4$ & $\sim 4$ \\
$10^5$ & $\sim 10^6$ & $\sim 6$ \\
$10^7$ & $\sim 10^8$ & $\sim 8$ \\
$10^{10}$ & $\sim 10^{11}$ & $\sim 11$ \\
\bottomrule
\end{tabular}
\end{center}
At $t > 10^7$, fewer than $8$ significant digits remain in the trigonometric evaluation,
degrading zero detection to tolerances worse than $10^{-8}$. At $t > 10^{10}$, the
argument reduction error ($\sim 10^{-5}$) exceeds the zero-detection threshold, rendering
double-precision computation unreliable.

For scientifically rigorous computation at large heights ($t > 10^7$), arbitrary-precision
arithmetic (e.g., MPFR~/ GNU~MP, or at minimum IEEE~754 quad precision via
\texttt{\_\_float128} or \texttt{long double} on x86) is required. The Riemann--Siegel
formula (Section~\ref{sec:rs}) partially mitigates this issue by requiring only
$M = O(\sqrt{t})$ terms (reducing the maximum argument to
$t\ln M \sim t\ln\sqrt{t} = (t/2)\ln t$), but the theta function evaluation
$\vartheta(t)$ itself involves arguments of order~$t$, so the precision wall persists.
Our current implementation is therefore limited to $t < 10^7$ for reliable zero detection
in double precision, with reduced reliability up to $t \sim 10^{10}$.

\paragraph{Cancellation near zeros.}
When $L(\beta, t)$ is near zero, the partial sum involves extensive cancellation. The
absolute value of~$L$ near a zero is $O(N^{-\beta})$, while individual terms have
magnitude $O(n^{-\beta})$ with the largest terms $O(1)$. Without compensated summation,
the absolute error in the sum is $O(N\,\eps_{\mathrm{mach}})$ (each of $N$ additions
contributes $O(\eps_{\mathrm{mach}})$ rounding error). With Kahan summation, the absolute
error is reduced to $O(\eps_{\mathrm{mach}})$, independent of~$N$. The value of~$L$ near
a zero is $O(N^{-\beta})$. The relative precision is therefore:
\begin{align}
  \text{Without Kahan:}\quad &\text{relative error}
    \sim \frac{N\,\eps_{\mathrm{mach}}}{N^{-\beta}}
    = N^{1+\beta}\,\eps_{\mathrm{mach}}, \label{eq:no-kahan}\\
  \text{With Kahan:}\quad &\text{relative error}
    \sim \frac{\eps_{\mathrm{mach}}}{N^{-\beta}}
    = N^{\beta}\,\eps_{\mathrm{mach}}. \label{eq:with-kahan}
\end{align}
For $N = 10^4$ and $\beta = 0.5$: without Kahan, relative precision
$\sim 10^6 \times 10^{-16} = 10^{-10}$. With Kahan summation:
$\sim N^\beta\,\eps_{\mathrm{mach}} = 10^2 \times 10^{-16} = 10^{-14}$.

This is adequate for zero detection with tolerance $10^{-6}$.


\subsection{Underflow Guard for $\log|L_N|$}
\label{sec:underflow}

The rate function $F_1(\beta, t) = -(1/\ln N)\ln|L_N(\beta, t)|$ involves taking the
logarithm of $|L_N|$. In floating-point arithmetic, $|L_N|$ can underflow to exactly
zero (when $\beta$ is large or at a near-zero of $\eta$ with large~$N$), producing
$\log(0) = -\infty$ and corrupting downstream computations.

\paragraph{Guard implementation.}
We use:
\begin{equation}\label{eq:underflow-guard}
  F_1 = -\frac{1}{\ln N}\,\ln\bigl(\max\bigl(|L_N|,\, 10^{-300}\bigr)\bigr).
\end{equation}
The floor value $10^{-300}$ is well above the IEEE~754 double-precision minimum subnormal
($\sim 5 \times 10^{-324}$) but far below any physically meaningful value. When
$|L_N| < 10^{-300}$, the computed $F_1$ is flagged as a sentinel value indicating
numerical underflow, and the point is marked for investigation with higher precision or
larger~$N$.

Additionally, any $(\beta, t)$ point where $|L_N| = 0$ exactly (to double precision) is
recorded separately and not included in the rate function analysis, since it may
represent either a genuine zero (signaling a DQPT) or a floating-point artifact.


\subsection{Grid Discretization Error}
\label{sec:grid-error}

Zeros detected by bisection on a grid with spacing $\Delta t$ have position error bounded
by $\Delta t/2$ in the initial detection, refined to arbitrary precision by iterative
bisection. After $k$ bisection iterations, the position error is
$\Delta t \cdot 2^{-(k+1)}$. For $\Delta t = 0.1$ and $k = 30$ iterations: position
error $\sim 10^{-10}$.

In the $\beta$ direction, the grid spacing $\Delta\beta$ determines our ability to
localize off-critical-line zeros. A zero at $\beta_0$ will be detected (in the sense
that $|L|$ dips below~$\eps$) only if $\beta_0$ is within $\sim\!\Delta\beta$ of a grid
point. For $\Delta\beta = 0.05$, we could miss a zero with $\beta_0 = 0.52$ (only $0.02$
from the critical line). However, the minimum of $|L(\beta, t)|$ over~$t$ is a smooth
function of~$\beta$, so a zero at $\beta_0 = 0.52$ would also produce a small value at
the nearest grid point $\beta = 0.50$ (the critical line)---making it detectable, though
its $\beta$-position would initially be misattributed to~$0.50$. The refinement stage
(Stage~4) would then resolve its true position.


\subsection{False Positive and False Negative Rates}
\label{sec:false-rates}

\paragraph{False positives} (declaring a zero where none exists) arise from numerical
artifacts such as:
\begin{itemize}
  \item Two nearly-cancelling large terms producing a spuriously small sum.
  \item Floating-point underflow near a local minimum of $|L|$.
\end{itemize}
Mitigation: every detected zero is confirmed by (a)~increasing~$N$, (b)~computing both
$L$ and~$G$, and (c)~verifying consistency with the Riemann--von Mangoldt zero count.

\paragraph{False negatives} (missing a genuine zero) arise from:
\begin{itemize}
  \item Grid spacing too coarse relative to the zero separation.
  \item Insufficient $N$ causing the partial sum to not resolve the zero.
\end{itemize}
Mitigation: grid spacing is chosen to be at most $1/4$ of the expected zero separation
(Section~\ref{sec:bisection}). The zero count check (Section~\ref{sec:controls},
Control~3) provides an independent confirmation that no zeros were missed.
